%% ============================================================================
%% Section: Introduction
%% ============================================================================
\section{Introduction}\label{sec:intro}

Causal directed acyclic graphs (DAGs) are widely adopted as the reference
model for reasoning about cause and effect across epidemiology, public health,
economics, and the empirical social sciences, to name a few. In a causal DAG
the nodes are variables and the edges encode direct causal
influence~\cite{pearl2009causality}, and analysts interrogate the graph with
\textit{structural} questions: which variables confound the effect of a
treatment $T$ on an outcome $Y$, whether a valid adjustment set exists and
which one it is, or whether $X$ is independent of $Y$ once $Z$ is held fixed.
Each of these questions is answered from the topology of the DAG alone. Such
DAGs are no longer small hand-drawn diagrams: they are large, continually
revised, and increasingly stored and versioned inside graph databases next to
the domain data they describe~\cite{pachera2025whatif,qiu2026causal}.

However, causal analysis is today divorced from the systems that store the
graph. The prevailing practice is to export the DAG into a standalone library
(e.g., DoWhy~\cite{sharma2020dowhy}, dagitty~\cite{textor2016dagitty},
pgmpy~\cite{ankan2024pgmpy}, or causal-learn~\cite{zheng2024causallearn}),
compute the answer in a separate process, and discard the in-memory
representation afterwards. This round-trip forfeits precisely the properties a
database is built to provide: durable and concurrent access to a single
authoritative model, declarative querying, integration with co-located data,
and incremental response to a changing graph. When an edge is added or a
variable is reclassified as latent, the whole model has to be re-extracted and
every downstream structure rebuilt from scratch. To the best of our knowledge,
no graph DBMS treats causal structural queries as a first-class, indexed
workload.

Previous work touching this problem follows three lines. (i) Causal
libraries~\cite{sharma2020dowhy,textor2016dagitty,ankan2024pgmpy,zheng2024causallearn,lee2023ananke}
implement d-separation, the adjustment criterion, and identification over an
in-memory graph, with efficient algorithms for the identification
problems~\cite{vanderzander2019separators,vanderzander2020minimal}; none is
DBMS-resident. (ii) Database research on causality has addressed causal
inference over relational data, hypothetical what-if queries, and causal data
integration~\cite{salimi2018hypdb,salimi2020carl,galhotra2022hyper,youngmann2023integration},
where the causal model is an input to query answering rather than a stored
object with its own operators. (iii) Graph database research offers mature
reachability indexes~\cite{yildirim2010grail,cheng2013tflabel,zhou2017dagreduction}
and declarative pattern languages~\cite{deutsch2022gql}, but neither
d-separation nor adjustment-set enumeration is a path pattern, and neither has
been treated as an indexed workload. These lines face a common limitation:
the operators the analyst needs are implemented far from the store, and the
store offers no operator that understands edge orientation and conditioning.

\begin{figure*}[t]
    \centering
    %% ============================================================================
%% Figure: running example -- fragment of the COVID-19 causal DAG
%% (Steiger et al. 2021). Edges are faithful to dev/seed-dag.cypher.
%% Included from sections/01_introduction.tex inside a figure* env.
%% Requires: \usepackage{tikz}; \usetikzlibrary{arrows.meta,positioning}
%% ============================================================================
\begin{tikzpicture}[
    font=\scriptsize,
    >={Stealth[length=1.8mm]},
    node distance=6mm,
    obs/.style={draw, rounded corners=1.5pt, minimum height=5.5mm,
                inner sep=2.5pt, align=center, fill=white},
    lat/.style={draw, dashed, rounded corners=1.5pt, minimum height=5.5mm,
                inner sep=2.5pt, align=center, fill=white},
    tnode/.style={obs, fill=black!12, very thick},
    ynode/.style={obs, fill=black!12, very thick},
    conf/.style={obs, draw=black, densely dotted, very thick},
    edge/.style={->, gray!70},
    bd/.style={->, thick},               % backdoor-path edge
    every node/.style={outer sep=0pt}
]

% --- nodes (placed by topological flow, left -> right) ---
\node[conf]  (burden) at (0.0, 1.5)  {covid19\_burden \\ \textit{(confounder)}};
\node[tnode] (interv) at (0.0, 0.0)  {interventions \\ ($T$)};
\node[obs]   (mob)    at (2.6, 0.8)  {mobility};
\node[lat]   (exp)    at (5.0, 0.4)  {exposure\_sarscov2};
\node[obs]   (acc)    at (6.6, 1.7)  {access\_tests};
\node[obs]   (age)    at (7.8, -1.2) {age};
\node[obs]   (sus)    at (10.2, -1.2){susceptibility};
\node[lat]   (onset)  at (7.8, 0.4)  {onset\_covid19};
\node[ynode] (rep)    at (10.2, 0.4) {reported\_cases \\ ($Y$)};

% --- causal edges ---
\draw[edge] (interv) -- (mob);
\draw[edge] (burden) -- (mob);
\draw[edge] (mob)    -- (exp);
\draw[edge] (exp)    -- (acc);
\draw[edge] (acc)    -- (onset);
\draw[edge] (age)    -- (onset);
\draw[edge] (age)    to[bend left=10] (exp);
\draw[edge] (sus)    -- (onset);

% --- backdoor path: T <- burden -> exposure -> onset -> reported ---
\draw[bd] (burden) -- (interv);
\draw[bd] (burden) to[bend left=12] (exp);
\draw[edge] (interv) to[bend right=12] (exp);
\draw[bd] (exp)    -- (onset);
\draw[bd] (onset)  -- (rep);

\end{tikzpicture}

    \caption{A fragment of the COVID-19 causal DAG of Steiger et
    al.~\cite{steiger2021covid} used as running example. Solid boxes are
    observed variables, dashed boxes are latent ones. \textit{interventions}
    ($T$) and \textit{reported\_cases} ($Y$) are the treatment and the outcome
    of the queries of Section~\ref{sec:intro}; \textit{covid19\_burden} is a
    common cause of both and opens the backdoor path drawn in bold.}
    \label{fig:example}
\end{figure*}

Fig.~\ref{fig:example} shows a fragment of a COVID-19 causal DAG curated by
epidemiologists~\cite{steiger2021covid}, stored as a property graph, on
which we consider three situations. (a) An analyst asks for the confounders
of \textit{interventions} and \textit{reported\_cases}: the answer,
\textit{covid19\_burden}, is the intersection of two ancestor sets, a query a
graph database can answer only with two variable-length path traversals
followed by a client-side intersection. (b) The analyst asks whether
\textit{susceptibility} is independent of \textit{age} given
\textit{reported\_cases}: the two variables share no ancestor and are
independent, but conditioning on \textit{reported\_cases}, a descendant of the
collider \textit{onset\_covid19}, opens the path between them, so the answer
depends on edge orientation and on the conditioning set, which no path pattern
expresses. (c) A curator adds an edge from \textit{access\_tests} to
\textit{reported\_cases} as new evidence accrues: every ancestor set computed
so far for \textit{reported\_cases} is stale, and an external tool has no way
to know it.

Based on the above observations, one can notice a gap between what the store
offers and what the workload needs. Three difficulties stand in the way of
closing it, and none is solved by wrapping a library behind a stored
procedure. First, the operators are not all plain reachability: d-separation
and adjustment-set construction depend on the conditioning set and on edge
orientation, and require a conditioning-aware traversal rather than a
transitive closure. Second, answered by traversing the stored graph on every
call, each query costs $O(|V|+|E|)$, and identification issues many such
queries, so per-query traversal does not scale to large DAGs or high query
volume. Third, the structure that accelerates these queries must survive
updates: a causal graph under active curation cannot afford a full
recompilation after every edit.

A key requirement for any solution is therefore a representation of the DAG
from which whole ancestor and descendant sets, not only point-to-point
reachability, can be read in constant or near-constant time, because every
operator in the causal catalog (confounders, mediators, the forbidden set of
an adjustment problem, the canonical adjustment set) is an algebraic
combination of such sets. Precomputed reachability closures stored as bitsets
meet this requirement: a closure is one array read, an intersection is one
word-parallel operation, and the linear-time identification algorithms of
van der Zander et al.~\cite{vanderzander2019separators} consume exactly these
sets. We therefore build the index around them and design the query
processing and the maintenance around what they make cheap.

We present \textbf{Causal4j}, a system that makes causal structural queries a
native, indexed workload inside a property-graph DBMS (Neo4j). Causal4j has
three parts, shown in Fig.~\ref{fig:arch}: a declarative causal query language
whose statements compile to physical operators over the graph store, a
compiled causal index that snapshots the stored DAG into a dense layout with
precomputed ancestor and descendant bitsets, and an incremental maintenance
scheme that keeps the index consistent as the graph is edited. In sum, we make
the following contributions.
\begin{itemize}
  \item We propose a declarative causal query language embedded in a graph
  DBMS, covering reachability, path roles, independence, identification, and a
  quantitative layer, together with its compilation from language operators to
  Cypher writes and physical procedures (Section~\ref{sec:lang}).
  \item We introduce a compiled reachability index for causal DAGs (dense CSR
  adjacency, topological layers, and ancestor/descendant bitsets) and the
  query-processing algorithms built on it: bitset algebra for reachability and
  path roles, ancestral masking for d-separation, and identification by
  covariate adjustment through the constructive criterion and polynomial-delay
  enumeration of minimal adjustment sets (Section~\ref{sec:index}).
  \item We formalize incremental maintenance of the index under edge and
  variable mutations, with an insertion patch and a deletion recompute both
  bounded by the affected cone of the graph, and we use the same machinery to
  derive interventional worlds by graph surgery (Section~\ref{sec:incr}).
  \item We evaluate Causal4j on standard Bayesian-network DAGs and synthetic
  DAGs against in-DBMS traversal, Cypher, Neo4j GDS, and external causal
  libraries. The results show that \TODO{headline: $N\times$ lower latency on
  identification, index build under $X$~ms, incremental update $Y$ against
  $Z$ for recompilation, exact agreement with dagitty on every verdict}
  (Section~\ref{sec:eval}).
\end{itemize}

\begin{figure}[t]
    \centering
    \resizebox{\linewidth}{!}{%% ============================================================================
%% Figure: Causal4j architecture / query pipeline
%% Included from sections/introduction.tex (or main.tex) inside a figure env.
%% Requires in main.tex preamble:
%%   \usepackage{tikz}
%%   \usetikzlibrary{arrows.meta,positioning,fit,backgrounds}
%% Single-column figure. For a wide (two-column) version, wrap the caller in
%% figure* and bump node distances.
%% ============================================================================
\begin{tikzpicture}[
    font=\footnotesize,
    >={Stealth[length=2mm]},
    box/.style={draw, rounded corners=2pt, align=center,
                minimum height=7mm, inner sep=3pt, fill=white},
    store/.style={draw, cylinder, shape border rotate=90, aspect=0.25,
                  align=center, minimum height=11mm, minimum width=15mm, fill=white},
    grp/.style={draw, dashed, rounded corners=3pt, inner sep=5pt},
    flow/.style={->, thick},
    node distance=5mm and 8mm
]

% --- language front-end ---
\node[box] (query) {Causal4j query \\ \texttt{\scriptsize IDENTIFY EFFECT ...}};
\node[box, right=of query] (ast) {Parser \\ + AST};
\node[box, right=of ast]  (trans) {Cypher \\ Translator};

\draw[flow] (query) -- (ast);
\draw[flow] (ast)   -- (trans);

% --- DBMS side ---
\node[box, below=10mm of trans, xshift=-4mm] (proc) {\texttt{CALL causal.*} \\ procedures};
\node[box, left=of proc] (writes) {Cypher \\ writes};

\draw[flow] (trans.south) .. controls +(0,-4mm) and +(0,4mm) .. (proc.north)
    node[midway, right] {reads};
\draw[flow] (trans.south) .. controls +(-8mm,-5mm) and +(0,4mm) .. (writes.north)
    node[midway, left] {schema};

% --- runtime / index ---
\node[box, below=9mm of proc, align=center] (index)
    {\textbf{Compiled causal index} \\[1pt]
     \scriptsize CSR adjacency $\cdot$ topo order/layer \\[-1pt]
     \scriptsize ancestor / descendant bitsets};
\node[store, left=30mm of index] (neo)
    {Neo4j \\ property \\ graph};

\draw[flow] (proc) -- (index) node[midway, right] {query};
\draw[flow] (writes.south) |- (neo.north) node[pos=0.75, above] {};
\draw[flow] ([yshift= 2mm]neo.east) -- ([yshift= 2mm]index.west)
    node[midway, above, font=\scriptsize] {compile};
\draw[flow] ([yshift=-2mm]index.west) -- ([yshift=-2mm]neo.east)
    node[midway, below, font=\scriptsize] {updates};

% --- results back to caller ---
\draw[flow] (index.east) .. controls +(14mm,0) and +(14mm,0) .. (proc.east);
\draw[flow] (proc.north east) .. controls +(10mm,10mm) and +(6mm,-6mm) .. (trans.south east)
    node[pos=0.5, right] {\scriptsize result rows};

% --- grouping boxes ---
\begin{scope}[on background layer]
  \node[grp, fit=(query)(ast)(trans), label={[font=\scriptsize\itshape]above:Language layer (client)}] {};
  \node[grp, fit=(writes)(proc)(index)(neo),
        label={[font=\scriptsize\itshape]below:Graph DBMS (Neo4j + Causal4j plugin)}] {};
\end{scope}

\end{tikzpicture}}
    \caption{Causal4j architecture. A statement is parsed and compiled into
    Cypher writes (definitions) or calls to \texttt{causal.*} procedures
    (queries); the procedures read a compiled snapshot of the stored DAG held
    by the plugin, which a transaction listener keeps consistent under
    updates.}
    \label{fig:arch}
\end{figure}

The remainder of the paper is structured as follows.
Section~\ref{sec:bg} defines causal DAGs, the query catalog, and the problem.
Section~\ref{sec:lang} presents the language and its compilation, and
Section~\ref{sec:compile} the compilation of the stored graph into a snapshot.
Section~\ref{sec:index} presents the index and the query-processing
algorithms, Section~\ref{sec:incr} incremental maintenance, and
Section~\ref{sec:quant} the quantitative layer. Section~\ref{sec:eval} reports
the experimental evaluation and Section~\ref{sec:related} discusses related
work. Finally, Section~\ref{sec:conclusion} concludes the paper.
